\setClassNameDisplay{LehrkraftKlasse}{lehrkraft\_unterrichtet\_klasse}

\ifbeamer
\Hefteintrag{1}{Darstellung von m:n-Beziehungen}{
m:n-Beziehungen können im UML-Klassendiagramm auf zwei verschiedene Arten dargestellt werden:
\vspace{0.3cm}

\emphColA{1. \LoesungLuecke{als direkte Beziehung}{15cm}}\\
Vorteil: \LoesungLuecke{Diagramm kompakt und übersichtlich}{15cm}

\vspace{0.3cm}
\begin{center}
\begin{tikzpicture}[scale=0.85, transform shape]
    \begin{class}[text width=3.5cm]{Lehrkraft}{0,0}
        \attribute{int id}
        \attribute{String name}
    \end{class}
    \begin{class}[text width=3.5cm]{Schulklasse}{{\linewidth/0.85-3.9cm},0}
        \attribute{int id}
        \attribute{String name}
    \end{class}
    \Loesung{
        \biAssociationSimple{Lehrkraft.east}{m}{lehrkraft\_unterrichtet\_klasse}{n}{Schulklasse.west}
    }
\end{tikzpicture}
\end{center}
}

\UnterHefteintrag{1}{Darstellung von m:n-Beziehungen}{
\emphColA{2. \LoesungLuecke{mit Beziehungstabelle}{15cm}}\\
Vorteil: \LoesungLuecke{direkt in relationaler Datenbank umsetzbar}{15cm}

\vspace{0.3cm}
\begin{center}
\begin{tikzpicture}[scale=0.82, transform shape]
    \begin{class}[text width=3.5cm]{Lehrkraft}{0,0}
        \attribute{int id}
        \attribute{String name}
    \end{class}
    \Loesung{
        \begin{class}[text width=6.5cm]{LehrkraftKlasse}{{(\linewidth/0.82-3.9cm)/2},0}
            \attribute{int lehrkraft\_id}
            \attribute{int schulklasse\_id}
        \end{class}
        \uniAssociation{[shift={(0,0.3)}]LehrkraftKlasse.west}{n}{lehrkraft\_id}{1}{[shift={(0,0.3)}]Lehrkraft.east}
        \uniAssociation{[shift={(0,0.3)}]LehrkraftKlasse.east}{n}{schulklasse\_id}{1}{[shift={(0,0.3)}]Schulklasse.west}
    }
    \begin{class}[text width=3.5cm]{Schulklasse}{{\linewidth/0.82-3.9cm},0}
        \attribute{int id}
        \attribute{String name}
    \end{class}
\end{tikzpicture}
\end{center}
}
\else
\Hefteintrag{1.5}{Darstellung von m:n-Beziehungen}{

m:n-Beziehungen können im UML-Klassendiagramm auf zwei verschiedene Arten dargestellt werden:
\vspace{0.3cm}

\emphColA{1. \LoesungLuecke{als direkte Beziehung}{15cm}}\\
Vorteil: \LoesungLuecke{Diagramm kompakt und übersichtlich}{15cm}

\vspace{0.3cm}
\begin{center}
\begin{tikzpicture}[scale=0.75, transform shape]
    \begin{class}[text width=3.5cm]{Lehrkraft}{0,0}
        \attribute{int id}
        \attribute{String name}
    \end{class}
    \begin{class}[text width=3.5cm]{Schulklasse}{{\linewidth/0.75-3.9cm},0}
        \attribute{int id}
        \attribute{String name}
    \end{class}
    \Loesung{
        \biAssociationSimple{Lehrkraft.east}{m}{lehrkraft\_unterrichtet\_klasse}{n}{Schulklasse.west}
    }
\end{tikzpicture}
\end{center}

\vspace{0.5cm}

\emphColA{2. \LoesungLuecke{mit Beziehungstabelle}{15cm}}\\
Vorteil: \LoesungLuecke{direkt in relationaler Datenbank umsetzbar}{15cm}

\vspace{0.3cm}
\begin{center}
\begin{tikzpicture}[scale=0.75, transform shape]
    \begin{class}[text width=3.5cm]{Lehrkraft}{0,0}
        \attribute{int id}
        \attribute{String name}
    \end{class}
    \Loesung{
        \begin{class}[text width=7.5cm]{LehrkraftKlasse}{{(\linewidth/0.75-3.9cm)/2},0}
            \attribute{}
        \end{class}
        \uniAssociation{[shift={(0,0.3)}]LehrkraftKlasse.west}{n}{lehrkraft\_id}{1}{[shift={(0,0.3)}]Lehrkraft.east}
        \uniAssociation{[shift={(0,0.3)}]LehrkraftKlasse.east}{n}{schulklasse\_id}{1}{[shift={(0,0.3)}]Schulklasse.west}
    }
    \begin{class}[text width=3.5cm]{Schulklasse}{{\linewidth/0.75-3.9cm},0}
        \attribute{int id}
        \attribute{String name}
    \end{class}
\end{tikzpicture}
\end{center}

}
\fi